% ===========================================================================
\chapter{The command line}\label{ch:commandline}

The last program puts it all together. \code{nex datalog.nex social.dl}
reads the file, checks it, evaluates it, answers the questions it asks,
and then waits for more questions, until the input ends or you type
\code{:quit}. Given a question on the command line as well, it answers that
one and stops.

\begin{newfile}{datalog.nex}
-- datalog.nex: run a Datalog program, then answer questions about it.
--
--   nex datalog.nex social.dl                     -- then ask questions
--   nex datalog.nex social.dl "sees(ann, X)."     -- ask one, and stop
--
-- A question is an atom, with or without "?-" before it. "why" before a
-- fact explains it; ":facts" lists every fact, ":stats" says how much
-- work the evaluation did, and ":quit" (or the end of the input) stops.

intern Parser
intern Analysis
intern Query
intern Explain
intern io/Path

-- The full path of a file the user named. The `nex` launcher runs
-- programs from its own installation folder, so a relative name is
-- resolved against the folder the program was started from.
function user_path(raw: String): String
do
  if raw.starts_with("/") then
    result := raw
  else
    result := program.working_directory + "/" + raw
  end
end
\end{newfile}

\code{user_path} is needed because the \code{nex} launcher runs programs
from its own folder, so a file name the user typed has to be found relative
to the folder they ran it in, which Nex calls \code{program.working_directory}.

\section{Reading a question}

A question can be typed with or without \code{?-} and the full stop. The
easiest way to read one is to make it into a one-question program and let
the parser do the work:

\begin{continued}{datalog.nex}
-- The one atom in `text`, a question or a fact, or why there is none.
function read_atom(text: String): Result[Atom, String]
do
  let t := text.trim
  if t.starts_with("?-") then
    t := t.substring(2, t.length).trim
  end
  if not t.ends_with(".") then
    t := t + "."
  end
  match parse_program("?- " + t) of
    Ok as ok then result := create Ok[Atom, String].make(ok.value.queries.get(0))
    Err as e then result := create Err[Atom, String].make(e.error)
  end
end
\end{continued}

\section{A session}

A \code{Session} holds an evaluated program: an explained database, so that
\code{why} works, and the semi-naive evaluator, whose counters
\code{:stats} reports. \code{respond} turns one line of input into the lines
to print.

\begin{continued}{datalog.nex}
class Session
create
  make(p: Datalog_Program)
  do
    source := p
    db := create Explained_Database.make
    evaluator := create Semi_Naive_Evaluator.make
    evaluate_program(p, db, evaluator)
  end
feature
  source: Datalog_Program
  db: Explained_Database
  evaluator: Semi_Naive_Evaluator

  summary: String
  do
    let given := 0
    across db.predicates as p do
      across db.facts_of(p) as f do
        if not db.is_derived(f) then
          given := given + 1
        end
      end
    end
    result := given.to_string + " facts given, " + (db.count - given).to_string
              + " derived in " + evaluator.rounds.to_string + " rounds."
  end
\end{continued}

\begin{continued}{datalog.nex}
  -- What to print in answer to one line of input.
  respond(line: String): Array[String]
  do
    let t := line.trim
    if t = ":stats" then
      result := [db.count.to_string + " facts; the evaluation took "
                 + evaluator.rounds.to_string + " rounds and looked at "
                 + evaluator.attempts.to_string + " facts."]
    elseif t = ":facts" then
      result := []
      across db.predicates as p do
        across db.facts_of(p) as f do
          result.add(f.text + ".")
        end
      end
    elseif t.starts_with("why ") then
      result := explain(t.substring(4, t.length))
    else
      match read_atom(t) of
        Ok as q then result := answer_lines(q.value, db)
        Err as e then result := [e.error]
      end
    end
  end

  explain(text: String): Array[String]
  do
    match read_atom(text) of
      Ok as a then
        if not a.value.is_ground then
          result := ["why needs a fact, without variables"]
        elseif not db.contains(a.value) then
          result := [a.value.text + " does not hold"]
        else
          result := db.why(a.value)
        end
      Err as e then result := [e.error]
    end
  end
end
\end{continued}

\section{The program}

\code{main} reports every kind of mistake in the input as a message: no
file, a file that does not parse, a program that is unsafe or cannot be
stratified. Only a program that passes all of those is evaluated, so the
preconditions of \code{evaluate_program} and the evaluators are never broken
by bad input; they are there to catch bugs in the engine.

\begin{continued}{datalog.nex}
function print_lines(lines: Array[String])
do
  across lines as l do
    console.print_line(l)
  end
end

function main()
do
  let args := program.arguments
  if args.is_empty then
    console.print_line("usage: nex datalog.nex PROGRAM.dl [QUESTION]")
  else
    let file := create Path.make(user_path(args.get(0)))
    if not file.exists then
      console.print_line("no such file: " + args.get(0))
    else
      match parse_program(file.read_text) of
        Err as e then console.print_line(args.get(0) + ", " + e.error)
        Ok as ok then
          if program_problem(ok.value) /= "" then
            console.print_line(args.get(0) + ": " + program_problem(ok.value))
          else
            converse(create Session.make(ok.value), args.get(0), args)
          end
      end
    end
  end
end

function converse(s: Session, name: String, args: Array[String])
do
  console.print_line(name + ": " + s.summary)
  across s.source.queries as q do
    console.print_line("?- " + q.text + ".")
    print_lines(answer_lines(q, s.db).map(fn (l: String): String do
      result := "   " + l
    end))
  end
  if args.length > 1 then
    print_lines(s.respond(args.get(1)))
  else
    let done := false
    from until done do
      console.print("?- ")
      let line := console.read_line
      if line = nil or line.trim = ":quit" then
        console.new_line
        done := true
      elseif line.trim /= "" then
        print_lines(s.respond(line))
      end
    end
  end
end
\end{continued}

\begin{continued}{datalog.nex}
main()
\end{continued}

\section{Running it}

\runoutput{datalog.nex social.dl "friends(ann, Who)."}{social}

The two questions in the file are answered first, then the one on the
command line. Without that one, the program would go on to wait for
questions to be typed.

\runoutput{datalog.nex social.dl "why suggest(ann, cy)."}{why_suggest}

\runoutput{datalog.nex flights.dl "airport(A)."}{flights}

The flights program uses a test, \code{Km < 1000}, to keep only short hops,
and the anonymous variables \code{_} for the arguments \code{airport} does
not care about. From Lisbon, flying only short hops, Paris is out of reach:
the flight from Madrid to Paris is 1050 km.

\runoutput{datalog.nex access.dl "refused(ann, What)."}{access}

\runoutput{datalog.nex chain.dl ":stats"}{chain_stats}

\runoutput{datalog.nex nowhere.dl}{no_file}
